\begin{abstract}
회전 기반 KV 캐시 양자화 방법들(TurboQuant, KVTC, PolarQuant, SpinQuant 등)은 모두 직교군 $SO(d)$의 원소를 키 벡터에 적용하지만, RoPE 하에서 어떤 회전이 왜 타당한지는 충분히 정리되지 않았다. 본 논문은 Class~C (직교 회전 + 독립 스칼라 양자화) 안에서 Pre-RoPE 공분산의 PCA가 MSE 기준 최적 회전이 됨을 소스 분포 가정 없이 증명한다 (정리~\ref{thm:optimal_hamiltonian}). 나아가, trained transformer에서 key 공분산과 query 공분산의 고유 구조가 자연적으로 정렬되어 있음 (0.6--2.5°)을 발견하고, 이를 기반으로 PCA 회전은 유지하면서 비트 배분에만 query 정보를 반영하는 Query-Weighted Water-Filling (QW-WF)을 제안하여 3모델 2-bit에서 10--33\% PPL을 추가 개선한다. 4개 모델(Qwen2.5-7B, Llama-3.1-8B, Mistral-7B, Qwen2.5-14B)에서의 검증 결과, Pre-RoPE PCA의 Pre-vs-Post 우위는 3-bit에서 4/4 모델로 확인되고, KVTC의 공유 PCA 대비 헤드별 PCA는 46.3\% PPL을 개선한다. 한편, attention-weighted 최적 회전(QW-PCA)과 MSE 최적 양자화기(Lloyd-Max)는 PPL에서 실패하며, 이에 대한 5가설 진단을 통해 ``회전 선택은 증명으로 해결, 비트 배분은 query 정보 활용, 양자화기 설계는 open problem''이라는 Class~C의 구조적 경계를 밝힌다.
\end{abstract}
